\section{Comparison results}
\label{section:comparison_result}

In this section, we compare the finite-state devices constructed from regular
expressions by the implementation in the directory \texttt{impl} of the accompanying
repository. For each expression, we construct
\begin{itemize}
  \item a finite-state automaton (NFA) using the algorithm described in
    \Cref{thm:thompson},
  \item a deterministic finite-state automaton (DFA) using the algorithm described in
    \Cref{thm:rabin_scott}, restricted to the states reachable from the initial state,
  \item $2$-separated and $3$-separated quasi-deterministic finite-state automata
    (2-QDFA, 3-QDFA) using \Cref{alg:quasi_det_sfsa},
  \item $2$-separated and $3$-separated quasi-deterministic finite-state automata
    (2-QDFA'', 3-QDFA'') using \Cref{alg:quasi_det_sfsa_components}.
\end{itemize}
The separator set used for the quasi-deterministic finite-state automata is
constructed using \Cref{alg:thompson_separator}. The last expression is a larger
instance of the previous one.
\bigbreak
\Cref{tbl:state_comparison} presents the number of states. For NFA and DFA it is the
cardinality of the set $Q$. For the quasi-deterministic finite-state automata it is the
sum of the cardinalities of the sets of quasi-states of the components, including the
empty quasi-state of each component.

\begin{table}[ht]
  \centering
  \small
  \setlength{\tabcolsep}{4pt}
  \begin{tabular}{||c c c c c c c||}
    \hline
    Expression & NFA & DFA & 2-QDFA & 3-QDFA & 2-QDFA'' & 3-QDFA'' \\ [0.5ex]
    \hline
    $(a.b.c...)$ & 52 & 27 & 31 & 36 & 108 & 216 \\
    \hline
    $(a|b|c|...).(a|b|c|...)*$ & 206 & 53 & 57 & 62 & 85 & 142 \\
    \hline
    $(0|\epsilon).(1|\epsilon)...$ & 60 & 11 & 15 & 20 & 29 & 40 \\
    \hline
    $(0|1)*.0.(0|1)^{16}$ & 106 & 131073 & 12353 & 1488 & 24670 & 3442 \\
    \hline
    $(0|1)*.0.(0|1)^{20}$ & 130 & 2097153 & 49409 & 3232 & 98686 & 7282 \\
    \hline
  \end{tabular}
  \caption{Comparison of the number of states of finite-state devices constructed
    from a regular expression.}
  \label{tbl:state_comparison}
\end{table}

\Cref{tbl:transition_comparison} presents the number of transitions. For NFA and DFA
it is respectively the cardinality of the sets $\Delta$ and $\delta$. For the
quasi-deterministic finite-state automata it is the number of entries of the functions
$\tau_1, \ldots, \tau_k$ with nonempty values, plus the number of entries of $\eta_1,
\ldots, \eta_k$ and $\eta_x$.

\begin{table}[ht]
  \centering
  \small
  \setlength{\tabcolsep}{4pt}
  \begin{tabular}{||c c c c c c c||}
    \hline
    Expression & NFA & DFA & 2-QDFA & 3-QDFA & 2-QDFA'' & 3-QDFA'' \\ [0.5ex]
    \hline
    $(a.b.c...)$ & 51 & 26 & 57 & 63 & 278 & 588 \\
    \hline
    $(a|b|c|...).(a|b|c|...)*$ & 257 & 1378 & 784 & 529 & 1636 & 2180 \\
    \hline
    $(0|\epsilon).(1|\epsilon)...$ & 69 & 55 & 52 & 53 & 168 & 260 \\
    \hline
    $(0|1)*.0.(0|1)^{16}$ & 124 & 262146 & 24770 & 3644 & 82238 & 11506 \\
    \hline
    $(0|1)*.0.(0|1)^{20}$ & 152 & 4194306 & 99074 & 8060 & 328958 & 24306 \\
    \hline
  \end{tabular}
  \caption{Comparison of the number of transitions of finite-state devices
    constructed from a regular expression.}
  \label{tbl:transition_comparison}
\end{table}

\Cref{alg:quasi_det_sfsa_components} produces larger devices than
\Cref{alg:quasi_det_sfsa}, since it also tabulates quasi-states that are not
components of reachable states. Both are subject to the bound $2^{|Q_i| + |X|}$ per
component, and both are much smaller than the DFA for the last two expressions.
\bigbreak
\Cref{tbl:time_comparison} presents the construction effort. By
\Cref{prop:quasidet_iterations}, the main loop of \Cref{alg:quasi_det_sfsa} executes
exactly one iteration per state of the DFA. The main loop of
\Cref{alg:quasi_det_sfsa_components} executes $\sum_{i=1}^k |U_i|$ iterations
(\Cref{prop:quasidet_components_algorithm}). The running times were measured on a
single run on an Apple silicon laptop and are only indicative.

\begin{table}[ht]
  \centering
  \begin{tabular}{||c c c c c||}
    \hline
    Expression & 2-QDFA & 3-QDFA & 2-QDFA'' & 3-QDFA'' \\ [0.5ex]
    \hline
    $(a.b.c...)$ & 28 & 28 & 108 & 216 \\
    \hline
    $(a|b|c|...).(a|b|c|...)*$ & 53 & 53 & 85 & 142 \\
    \hline
    $(0|\epsilon).(1|\epsilon)...$ & 12 & 12 & 29 & 40 \\
    \hline
    $(0|1)*.0.(0|1)^{16}$ & 131073 & 131073 & 24669 & 3441 \\
     & 472 ms & 629 ms & 41 ms & 5 ms \\
    \hline
    $(0|1)*.0.(0|1)^{20}$ & 2097153 & 2097153 & 98685 & 7281 \\
     & 9 s & 11 s & 206 ms & 13 ms \\
    \hline
  \end{tabular}
  \caption{Number of iterations of the main loop and, for the larger expressions,
    construction time.}
  \label{tbl:time_comparison}
\end{table}

For the first three expressions both algorithms complete in a few milliseconds. For
the last two, the number of iterations of \Cref{alg:quasi_det_sfsa} grows with the
size of the DFA, whereas \Cref{alg:quasi_det_sfsa_components} only depends on the
sizes of the components and the separator. For $(0|1)*.0.(0|1)^{20}$ and $k = 3$, it
constructs the quasi-deterministic finite-state automaton more than $800$ times faster.
\bigbreak
\Cref{tbl:symbol_comparison} repeats the comparison of the number of states with the
symbol-balanced variant of \Cref{alg:thompson_separator}
(\Cref{thm:thompson_separator_symbols}), which balances the essential states instead of
all states. For $3$ components, the implementation applies either variant twice, as in
the tables above.

\begin{table}[ht]
  \centering
  \small
  \setlength{\tabcolsep}{4pt}
  \begin{tabular}{||c c c c c||}
    \hline
    Expression & 2-QDFA & 3-QDFA & 2-QDFA'' & 3-QDFA'' \\ [0.5ex]
    \hline
    $(a.b.c...)$ & 31 & 36 & 108 & 216 \\
    \hline
    $(a|b|c|...).(a|b|c|...)*$ & 57 & 62 & 164 & 366 \\
    \hline
    $(0|\epsilon).(1|\epsilon)...$ & 15 & 20 & 30 & 42 \\
    \hline
    $(0|1)*.0.(0|1)^{16}$ & 6273 & 976 & 12478 & 2098 \\
    \hline
    $(0|1)*.0.(0|1)^{20}$ & 49409 & 3232 & 98686 & 7282 \\
    \hline
  \end{tabular}
  \caption{Number of states of quasi-deterministic finite-state automata constructed
    with symbol-balanced separators.}
  \label{tbl:symbol_comparison}
\end{table}

For $(0|1)*.0.(0|1)^{16}$, the symbol-balanced separator yields components with $m_1 =
22$ and $m_2 = 14$ essential states instead of $24$ and $12$, and the
quasi-deterministic finite-state automata are about half as large. For $(0|1)*.0.(0|1)^{20}$
both variants find the same separator. For $(a|b|c|...).(a|b|c|...)*$ the
symbol-balanced separator is worse for \Cref{alg:quasi_det_sfsa_components}: both
separators split the essential states evenly, but the symbol-balanced one leads to more
quasi-states that are not components of reachable states. The bounds of
\Cref{lem:quasidet_essential} are worst-case bounds, and a separator that is better
with respect to them need not be better on a particular finite-state automaton.
\bigbreak
\Cref{tbl:boundary_comparison} presents the construction of
\Cref{section:quasi_determinization_boundaries} with the boundaries $X_i$ consisting of the
separator states that share a transition with $Q_i$, for up to $16$ components. The
separators are symbol-balanced and obtained as above. For each number of components, the
table shows the number of quasi-states $\sum_{i=1}^k |U_i^b \cup T_i^b|$. Since
$\eta_i^b(t, \emptyset) = t \cap \mathcal{B}(Q_i)$, we have $T_i^b \subseteq U_i^b$, so this is
also the number of iterations of the main loop. The last columns show $|X|$ and $\max_i |X_i|$ for $16$ components.

\begin{table}[ht]
  \centering
  \small
  \setlength{\tabcolsep}{4pt}
  \begin{tabular}{||c c c c c c c||}
    \hline
    Expression & $k = 2$ & $k = 4$ & $k = 8$ & $k = 16$ & $|X|$ & $\max_i |X_i|$ \\ [0.5ex]
    \hline
    $(a.b.c...)$ & 88 & 88 & 100 & 101 & 18 & 2 \\
    \hline
    $(a|b|c|...).(a|b|c|...)*$ & 110 & 162 & 240 & 299 & 30 & 4 \\
    \hline
    $(0|\epsilon).(1|\epsilon)...$ & 25 & 31 & 43 & 56 & 16 & 2 \\
    \hline
    $(0|1)*.0.(0|1)^{16}$ & 12477 & 421 & 177 & 137 & 22 & 2 \\
     & 22 ms & 669 $\mu$s & 219 $\mu$s & 154 $\mu$s & & \\
    \hline
    $(0|1)*.0.(0|1)^{20}$ & 98685 & 997 & 297 & 185 & 19 & 2 \\
     & 246 ms & 1 ms & 396 $\mu$s & 217 $\mu$s & & \\
    \hline
  \end{tabular}
  \caption{Number of quasi-states and, for the larger expressions, construction time of
    quasi-deterministic finite-state automata with boundaries.}
  \label{tbl:boundary_comparison}
\end{table}

For $(0|1)*.0.(0|1)^{20}$ with $16$ components, the quasi-deterministic finite-state
automaton has $185$ quasi-states, compared with $7282$ for $3$ components without boundaries
and $2097153$ states of the deterministic finite-state automaton. Without boundaries,
\Cref{alg:quasi_det_sfsa_components} would tabulate $\eta_i$ for $2^{19}$ separator states
per quasi-state and $\eta_x$ with $4^{19}$ entries. For $(a|b|c|...).(a|b|c|...)*$, whose
deterministic finite-state automaton has only $53$ states, more components do not help: the
boundaries have up to $4$ states, and the number of quasi-states grows with $k$. For the small expressions, the number of quasi-states grows slowly with $k$,
since every component contributes at least its start and empty quasi-states.
\bigbreak
\Cref{tbl:throughput_comparison} presents the matching throughput, measured with the program
\texttt{bench} of the implementation. All devices are converted to dense arrays, so that
every look-up in a table is a single array access. The NFA is simulated with bit sets: for
every state and symbol, the $\epsilon$-closure of the successors is precomputed, and a step
unites these sets for the active states. Each device matches a uniformly random word of
$4$ MiB over the alphabet of the expression in full. The table shows the median of $7$
runs on an Apple silicon laptop; repeated measurements differed by up to $15\%$. The DFA
for $(0|1)*.0.(0|1)^{24}$ with $2^{25} + 1$ states was not constructed.

\begin{table}[ht]
  \centering
  \small
  \setlength{\tabcolsep}{2pt}
  \begin{tabular}{||c c c c c c c c c||}
    \hline
    Expression & NFA & DFA & 2-QDFA'' & 3-QDFA'' & \multicolumn{4}{c||}{with boundaries} \\
    & & & & & $k = 2$ & $k = 4$ & $k = 8$ & $k = 16$ \\ [0.5ex]
    \hline
    $(a|b|c|...).(a|b|c|...)*$ & 8.9 & 355 & 106 & 84 & 142 & 72 & 58 & 23 \\
    \hline
    $(0|1)*.0.(0|1)^{16}$ & 15.9 & 192 & 65 & 83 & 88 & 112 & 74 & 29 \\
    \hline
    $(0|1)*.0.(0|1)^{20}$ & 11.8 & 69 & 52 & 61 & 62 & 111 & 74 & 29 \\
    \hline
    $(0|1)*.0.(0|1)^{24}$ & 9.9 & -- & 9.7 & 46 & 10.7 & 110 & 73 & 29 \\
    \hline
  \end{tabular}
  \caption{Matching throughput in MB/s.}
  \label{tbl:throughput_comparison}
\end{table}

The DFA is the fastest device as long as its table is small. For $(0|1)*.0.(0|1)^{20}$ its
table has $2^{22}$ entries, look-ups miss the cache, and its throughput drops to $69$ MB/s.
The quasi-deterministic finite-state automata show the same effect: with $2$ components,
the throughput drops to about $10$ MB/s for $(0|1)*.0.(0|1)^{24}$, where the components have
$787198$ quasi-states. With $4$ components and boundaries, the tables have at most $2005$
quasi-states, and the throughput stays at about $110$ MB/s for all three expressions. This
is faster than the DFA for $(0|1)*.0.(0|1)^{20}$ and about $11$ times faster than the
simulation of the NFA. Beyond that, the time per symbol grows linearly with the number of
components, as expected from \Cref{alg:qdfsa_traversal}, and the throughput drops to about
$29$ MB/s for $16$ components, which is still about three times the throughput of the NFA
simulation. For $(a|b|c|...).(a|b|c|...)*$, whose DFA has $53$ states, the DFA is best and
additional components only cost time. With $2$ components, the device with boundaries is
faster than the one of \Cref{alg:quasi_det_sfsa_components}, since it combines the
separator states by a bitwise union instead of look-ups in the table of $\eta_x$.

The number of components thus trades the size of the tables for the work per symbol. The
best choice in these experiments is the smallest number of components for which the tables
fit into the cache.
